\documentclass[../main.tex]{subfiles}
\begin{document}

\section{Malformed attack script deployment}
\label{app:malformed-script-deploy}

Chapter~\ref{sec:malformed_plc} describes four groups of malformed S7comm payloads and Table~\ref{tab:malformed-changes} lists the packet modifications used in the lab. This appendix maps each script \texttt{--mode} to those modifications and to Suricata SIDs, and records how to run the lab script.

\subsection{Script modes and Suricata SID mapping}

Each sample starts from a valid Setup Communication PDU (same hex as \texttt{s7comm\_dos.py}) and modifies only one field (or truncates) so results can be compared easily in Wireshark.

\begin{table}[H]
    \centering
    \footnotesize
    \begin{tabular}{|l|l|l|l|}
        \hline
        Group & \texttt{--mode} script & Change vs.\ valid Setup & SID \\
        \hline
        1 & \texttt{tpkt\_bad\_version} & TPKT byte 0 = \texttt{0x02} & 1000050 \\
        \hline
        1 & \texttt{tpkt\_len\_gt\_payload} & TPKT length = \texttt{0x0100}, body $\approx$ 36 B & 1000051 \\
        \hline
        1 & \texttt{tpkt\_len\_lt\_payload} & TPKT length = \texttt{0x0008}, full Setup body & 1000058 \\
        \hline
        1 & \texttt{cotp\_invalid} & Replace \texttt{02 F0 80} with \texttt{01 00 00} & 1000056 \\
        \hline
        2 & \texttt{s7\_bad\_protocol\_id} & Byte 7 = \texttt{0x31} & 1000052 \\
        \hline
        2 & \texttt{s7\_bad\_rosctr} & Byte 8 = \texttt{0xFF} & 1000053 \\
        \hline
        3 & \texttt{job\_bad\_opcode} & Byte 17 = \texttt{0xFF} (instead of \texttt{0xF0}) & 1000054 \\
        \hline
        3 & \texttt{write\_var\_malformed} & Job \texttt{0x05}, 20 B PDU (missing item) & 1000057 \\
        \hline
        4 & \texttt{setup\_truncated} & Send only the first 14 bytes & 1000055 \\
        \hline
    \end{tabular}
    \caption{Malformed script modes and detection rule mapping}
    \label{tab:malformed-phases}
\end{table}

\subsection{Running the lab script}

All scenarios are implemented in \texttt{attacks/malformed/s7comm\_malformed.py}. By default, \texttt{--mode all} runs 9 phases sequentially; each phase sends \texttt{--repeats} times (default 3) on a new TCP session to port 102, then waits \texttt{--pause} seconds (default 6) before the next phase---similar to DoS but without flooding: the goal is to produce clear malformed alerts in \texttt{fast.log}.

\begin{packetdump}
python3 attacks/malformed/s7comm_malformed.py <IP_PLC>
# equivalent to: --mode all --repeats 3 --pause 6
python3 attacks/malformed/s7comm_malformed.py <IP_PLC> --mode s7_bad_rosctr
\end{packetdump}

Processing on the wire: attacker $\rightarrow$ TCP/102 $\rightarrow$ mirror $\rightarrow$ Suricata. Unlike \texttt{ics\_dos.rules}, malformed rules do not use thresholds: each packet matching an anomalous signature may generate one alert (\texttt{sid:1000050}--\texttt{1000058}).

\end{document}
